\chapter{Outlook}
\label{chap:outlook}

This final chapter looks beyond the present work: at the immediate follow-up
studies motivated by the results, at the \PLATO mission timeline within which
those studies will operate, and at the complementary role of atmospheric
characterisation missions.


\section{Follow-up Publications}
\label{sec:follow_up}

\emph{Scaffold.}
Immediate follow-up directions. Publication of the \PLATO Input Catalog
application as a companion paper to \textcite{Schlecker2024}. Extension of
the framework to explicit H/He contamination modelling. Integration of the
\textcite{Aguichine2021} evolutionary tracks to replace the parametric
atmospheric-lifetime prescription used in this work.


\section{The PLATO Mission Timeline}
\label{sec:plato_timeline}

\emph{Scaffold.}
Expected timeline between mission launch, the completion of the first
long-duration observation phase, and the first \HZIED-capable data release.
Readiness of the analysis pipeline developed in this thesis for that data
release, and implications for the design of the analysis-team's early
science plan.


\section{Synergy with Atmospheric Characterisation}
\label{sec:synergy}

\emph{Scaffold.}
Complementarity between a statistical \HZIED detection by \PLATO and
individual-planet atmospheric confirmation with ELT/METIS, JWST, and future
missions such as LIFE and HWO. Path from a demographic detection to a
physical confirmation of the runaway greenhouse transition in a specific
target system.
